\subsection{Zero-addition gates: the minimum and an algorithm that attains it}\label{sec:zeroadd}

The ROT-preprocessing adders (\flag{ROT_PREPROCESSING_OPT}) use Triad Suite's function-dependent preprocessing: every
input of a multiplication gate takes a field of the gate's random Beaver tuple as its mask, so the tuple needs no
correction, and a wire whose mask would have to equal two different fields (fan-out into two multiplication gates, or
an XOR of masks that are fixed elsewhere) gets a \emph{zero-addition gate} on one of its edges: a new wire of value 0
whose mask, published in preprocessing, re-masks that edge. Each zero-add costs one opening in preprocessing
(\texttt{zero\_add}) and nothing online. Triad's Alg.~1/2 place them by propagating masks through the XOR gates and
resolving conflicts as they appear; this subsection gives the exact minimum, an algorithm that attains it by
construction, and the proof.

\paragraph{The linear model.} Let the circuit have XOR gates and multiplication gates (AND, AND3, AND4, dot terms; an
a-known product counts as a gate whose output carries a fresh mask, its remask or its dot chain's root). Call
\emph{sources} the masks of the circuit inputs that are not fixed (a fixed mask, such as the bake's prescribed masks
in the AB circuits, is a constant) and of the multiplication outputs, and let $s \in \mathbb{F}_2^n$ collect them.
Every wire's mask is then affine in $s$: an XOR gate adds its input masks. A \emph{constraint} is an edge into a
multiplication gate (the input pairs a reshared circuit reshares have none); with $m$ constraints, their masks are
$\Lambda s + d$ for an $m \times n$ matrix $\Lambda$ over $\mathbb{F}_2$ and a constant $d$. The tuple fields are
independent and uniform, so preprocessing must be able to make the constrained masks equal any vector $t$ of fields:
the affine map $s \mapsto \Lambda s + d$ must be onto, i.e.\ $\operatorname{rank}\Lambda = m$. A zero-add on
constraint $c$ replaces its mask by a new source of its own: it appends the unit column $e_c$ to $\Lambda$.

\textbf{Theorem 3.} The minimum number of zero-addition gates of a circuit is $m - \operatorname{rank}_{\mathbb{F}_2}\Lambda$.
A set $Z$ of constraints suffices if and only if the rows of $\Lambda$ outside $Z$ are linearly independent; with
such a $Z$ the masks of all opened wires are uniform and independent.

\emph{Proof.} \emph{Lower bound.} Appending a column raises the rank by at most one, so
$\operatorname{rank}[\Lambda \mid e_Z] \le \operatorname{rank}\Lambda + |Z|$, and full rank $m$ needs
$|Z| \ge m - \operatorname{rank}\Lambda$. \emph{Characterization.} If the rows outside $Z$ are independent, take a
linear combination of the rows of $[\Lambda \mid e_Z]$ that vanishes: each row $c \in Z$ is the only one with a 1 in
column $e_c$, so its coefficient is 0, and the remaining rows are independent, so all coefficients are 0; the rank is
$m$. Conversely, if the rows outside $Z$ are dependent, the same combination vanishes on $[\Lambda \mid e_Z]$ as well
(those rows are 0 in the columns $e_Z$), so the rank is below $m$. Rows outside a row basis of $\Lambda$ are such a set
of size $m - \operatorname{rank}\Lambda$, so the bound is attained. \emph{Masks.} Preprocessing draws the tuple fields
$t$ uniformly, chooses $s$ uniformly among the solutions of $\Lambda_{\bar Z} s + d_{\bar Z} = t_{\bar Z}$ (they exist
since $\Lambda_{\bar Z}$ has full row rank) and publishes $t_c - (\Lambda_c s + d_c)$ for $c \in Z$. As $t_{\bar Z}$ is
uniform and the map onto, $s$ is uniform on $\mathbb{F}_2^n$, so the masks of all multiplication outputs (the opened
wires) are jointly uniform; each published value contains its own fresh field $t_c$ and is uniform and independent of
the rest. This is the distribution of the protocol with fresh masks everywhere. \hfill$\square$

Over $\mathbb{Z}_{2^\ell}$ (arithmetic circuits) the masks are integer combinations with coefficients $\pm1$, and an
integer matrix is onto modulo $2^\ell$ exactly when some $m \times m$ minor is odd, i.e.\ when its reduction modulo 2
has rank $m$: Theorem~3 holds with the rank over $\mathbb{F}_2$ of $\Lambda \bmod 2$.

\begin{figure}[t]
\centering\small
\fbox{\begin{minipage}{0.95\linewidth}
\textbf{Algorithm RankFDP} (minimal zero-addition gates)\\[2pt]
\emph{Input:} a circuit; the fixed-mask inputs; the reshared gates; optionally a cost per constraint.
\emph{Output:} the zero-add edges $Z$ and every wire's mask.
\begin{enumerate}\itemsep1pt
\item \textbf{Forms.} In topological order: a non-fixed input or a multiplication output gets a new source (a unit
  bitset), a fixed input the constant 0, an XOR output the XOR of its inputs' forms. Each edge into a multiplication
  gate (not reshared) appends its form to the constraint list.
\item \textbf{Elimination.} Sort the constraints by decreasing cost (default: circuit order). Keep a basis indexed by
  pivot (the highest source of a row). For each constraint, reduce its form by the basis; if a non-zero remainder is
  left, it enters the basis with that pivot; otherwise put the edge into $Z$.
\item \textbf{Masks.} Draw the tuple fields; assign the sources that are pivots from the basis rows (each basis row
  fixes its pivot given the sources below it, solved from the lowest pivot up) and draw the others uniformly; every
  wire's mask is its form evaluated; an edge in $Z$ gets its tuple field and a zero-add publishing the difference.
\end{enumerate}
\end{minipage}}
\caption{RankFDP. With $n$ sources and $m$ constraints it takes $O(m \cdot \operatorname{rank}\Lambda \cdot n / w)$ word
operations ($w$-bit words), milliseconds for the 64-bit adders (\texttt{zero\_add\_bound.py} in the circuit generator).}
\label{fig:rankfdp}
\end{figure}

\paragraph{Optimality of RankFDP.} The sets of independent rows of $\Lambda$ form a matroid (linear independence), and
by Theorem~3 the complements of its bases are exactly the minimum zero-add sets. Step~2 builds a basis greedily, so
$|Z| = m - \operatorname{rank}\Lambda$ in every order. With costs it is the matroid greedy algorithm, which yields a
basis of maximum total cost~\cite{edmonds71}: its complement $Z$ has minimum total cost among all minimum zero-add sets
(and among all valid sets, since every valid set contains the complement of a basis). This matters where zero-adds
differ in price: under the reshare bake the zero-adds on input slices that P1 bakes are local
(\texttt{zero\_add\_local}), so they should be preferred to the ones that send their difference.

\paragraph{The generator is at the minimum.} \Cref{tab:zeroadd} compares the bound with the zero-adds the circuit
generator places with Alg.~1/2 for all twelve families the adders use, at 8, 16, 32 and 64 bits and the narrow widths:
it is at $m - \operatorname{rank}\Lambda$ everywhere, so these circuits cannot do with fewer zero-adds under
this model. (Its mask propagation never creates an avoidable dependency on them; RankFDP guarantees it for any
circuit.) The minimum is a property of the circuit's logic and of the tuple model, not of the placement. What remains
is the model: if the multiplication gates reading one wire shared that wire's tuple field (\emph{fan-out tuples}: one
constraint per wire instead of per edge), the minimum would fall by 21\% (PPA), 31\% (four-way PPA) and 70\% (a-known
four-way PPA, 104 to 31 at 32 bits), RCA unchanged. Such tuples are cheaper to generate as well (a shared factor's
share bits are the choice bits of one ROT with a multi-bit correlation instead of one ROT per product), but they need
a new tuple type in the generator, the OT preprocessing and the adder classes.

\begin{table}[t]
\centering\small
\begin{tabular}{@{}lrrrrrr@{}}
\toprule
circuit & $k$ & $m$ & rank & min & generator & fan-out min\\
\midrule
RCA MSB, AB & 32 / 64 & 62 / 126 & 30 / 62 & 32 / 64 & 32 / 64 & 32 / 64\\
RCA MSB, AB reshared & 32 / 64 & 60 / 124 & 30 / 62 & 30 / 62 & 30 / 62 & 30 / 62\\
RCA MSB, a-known & 32 / 64 & 60 / 124 & 29 / 61 & 31 / 63 & 31 / 63 & 31 / 63\\
PPA MSB, AB & 32 / 64 & 172 / 362 & 55 / 118 & 117 / 244 & 117 / 244 & 92 / 188\\
PPA MSB, AB reshared & 32 / 64 & 110 / 236 & 55 / 118 & 55 / 118 & 55 / 118 & 30 / 62\\
PPA MSB, a-known & 32 / 64 & 110 / 236 & 55 / 118 & 55 / 118 & 55 / 118 & 30 / 62\\
PPA4 MSB, AB & 32 / 64 & 154 / 331 & 21 / 46 & 133 / 285 & 133 / 285 & 92 / 188\\
PPA4 MSB, AB reshared & 32 / 64 & 132 / 289 & 21 / 46 & 111 / 243 & 111 / 243 & 70 / 146\\
PPA4 MSB, a-known & 32 / 64 & 125 / 271 & 21 / 46 & 104 / 225 & 104 / 225 & 31 / 63\\
PPA4 MSB, a-known (cut) & 27 / 52 & 105 / 220 & 18 / 37 & 87 / 183 & 87 / 183 & 26 / 51\\
\bottomrule
\end{tabular}
\caption{Zero-addition gates per adder: constraints $m$, $\operatorname{rank}\Lambda$, the minimum of Theorem~3, the
generator's count (Alg.~1/2) and the minimum with fan-out tuples (\texttt{zero\_add\_bound.py}; all widths 8--64 in
the script's output agree as well).}
\label{tab:zeroadd}
\end{table}
